\SourceLeaf{322}%
\par\noindent\hspace*{\parindent} Zunächst ergibt sich aus der Verbindung der Gleichungen (1), (2) und (3), dass \(qq+\alpha\alpha ll-uu-ss\) eine Größe vierter Ordnung ist. Daher kann man anstelle von (2) auch schreiben
\[
l=\frac{s\mathord{.}\sin U}{\alpha\mathord{.}\cos Q}(1-\tfrac1{24}qq+\tfrac1{24}uu)
\]
oder, wenn man nach (8), (9) und (10)
\[
s=\frac{r\cos\theta^2}{a\cos\varphi},\qquad \alpha\cos Q=\frac{\cos\theta\cos P}{\cos\varphi}
\]
schreibt, erhält man
\[
l=\frac{r\cos\theta\sin U}{a\cos P}(1-\tfrac1{24}qq+\tfrac1{24}uu)
\]
Mit Gleichung (4) und einer einfachen Rechnung lässt sich ferner \(\frac{\cos\theta}{\cos P}=\frac{\BareRoot{(1-ee\sin P^2)}}{\cos P}\) wie folgt entwickeln
\[
\frac{\cos\theta}{\cos P}=\frac{\BareRoot{(1-ee\sin B^2)}}{\cos B}(1+\frac{3ee\sin P^2}{8(1-ee\sin P^2)}\mathord{.}qq)
\]
Dieser Ausdruck ist bis zur vierten Ordnung ausschließlich genau. Schreibt man zugleich gemäß Gleichung (6) \(T\) für \(U\), so ergibt sich
\[
l=\frac{r\BareRoot{(1-ee\sin B^2)}\mathord{.}\sin T}{a\cos B}(1-\frac{1-10ee\sin P^2}{24(1-ee\sin P^2)}\mathord{.}qq+\tfrac1{24}uu)
\]
Im eingeklammerten Teil kann man noch \(\frac{\cos\varphi}{\BareRoot{(1-ee\sin P^2)}}\mathord{.}b\) für \(q\), dann \(t\) für \(u\) und schließlich \(B\) für \(P\) einsetzen. Wie man aus den Gleichungen (5), (7) und (4) leicht sieht, bleibt bei all diesen Ersetzungen der Grad der Genauigkeit erhalten. Schreiben wir außerdem zur Abkürzung
\[
\BareRoot{(1-ee\sin B^2)}=k
\]
so erhalten wir (I)
\[
l=\frac{kr\sin T}{a\cos B}(1-\frac{(1-10ee\sin B^2)\cos\varphi^2}{24k^4}\mathord{.}bb+\tfrac1{24}tt)
\]

\Article{24}
Auf ähnliche Weise geht Gleichung (1) über in
\[
q=s\cos U(1-\tfrac1{12}qq+\tfrac1{12}ss+\tfrac18uu)
\]
und damit Gleichung (5) in
